%% chapters/chapter05/chapter05.tex
%% CHAPTER 5: CONCLUSION

\chapterfigpath{chapter05}

\chapter{Conclusion}
\label{chap:conclusion}

\section{Summary of Research}
\label{sec:summary}

This research addresses catastrophic forgetting in edge robotics continual learning, where standard methods are either memory-heavy or computationally prohibitive. To resolve this, this thesis proposed \textbf{HippoCortex}, a bio-inspired, decoupled sleep-wake framework. HippoCortex combines a selective Mamba backbone \cite{guMambaLinearTimeSequence2024} with a Sleep-Wake Replay (SWR) CVAE generator and a Null-Space Gradient Projector.

The research was executed across several key phases:
\begin{enumerate}
    \item A systematic literature review analyzing the stability-plasticity dilemma and detailing the limitations of EWC \cite{kirkpatrickOvercomingCatastrophicForgetting2017}, ER \cite{riemer2019learning}, GPM \cite{sahaGradientProjectionMemory2021}, and Mamba-CL \cite{chengMambaCLOptimizingSelective2025}.
    \item Mathematical design of HippoCortex, including the stats buffer, CVAE-based hidden state generator, and SVD/QR null-space projection.
    \item Reproduction of the Mamba-CL \cite{chengMambaCLOptimizingSelective2025} baseline on 10-task Split-CIFAR100 \cite{venThreeScenariosContinual2019} to establish empirical targets.
    \item Architectural layout of the ROS2-based on-board deployment system for the Unitree Go2 quadruped.
\end{enumerate}

\section{Key Findings and Contributions}
\label{sec:contributions}

This research makes several key contributions to the field of continual learning and robotic edge intelligence:
\begin{itemize}
    \item \textbf{Decoupled Consolidation Paradigm:} Formulated a sleep-wake CL loop that decouples awake training (focused on immediate task plasticity) from sleep-phase consolidation (focused on stability protection). This design prevents on-board computational bottlenecks during the robot's active operation.
    \item \textbf{Raw-Data-Free Replay:} Designed a hidden-state statistics-based generative replay mechanism that scales as $\mathcal{O}(T \times d_{\text{model}})$ in memory complexity, compared to Experience Replay's memory scaling of $\mathcal{O}(N_{\text{samples}} \times D_{\text{raw}})$ where $D_{\text{raw}}$ is the raw input dimension. Storing only two vectors per task ($\mu_t, \sigma_t^2$) reduces the memory footprint to less than $20$\,KB, eliminating privacy risks and raw sensor storage constraints on edge platforms.
    \item \textbf{Rigorous Baseline Verification:} Successfully reproduced the Mamba-CL baseline on Split-CIFAR100, achieving a final average class-incremental accuracy of $90.16\%$ and restricting the forgetting measure to $2.18\%$. This confirms the efficacy of null-space gradient projection, but highlights a 17-hour training overhead that justifies the need for HippoCortex's decoupled consolidation.
\end{itemize}

The initial research objectives have been achieved as follows:
\begin{itemize}
    \item \textbf{Objective 1 (Literature Analysis):} Completed a comprehensive survey of 20+ sources, placing HippoCortex at the intersection of selective SSMs and generative replay.
    \item \textbf{Objective 2 (Methodological Formulation):} Mathematically formulated the CVAE loss, the stats buffer update, and the incremental projection update.
    \item \textbf{Objective 3 (Experimental Validation):} Conducted baseline evaluations, documenting task-by-task accuracy drift and computing latency profiles.
\end{itemize}

\section{Limitations}
\label{sec:limitations}

Despite the positive progress, several limitations of the current stage are acknowledged:
\begin{itemize}
    \item \textbf{Plasticity Decay Trade-off:} The baseline results confirm that as the null-space projection basis $U$ gains rank across the 10 tasks, the model's plasticity drops, leading to lower initial accuracies on later tasks (e.g., Task 10 accuracy starts at $89.60\%$ compared to Task 2 at $97.80\%$). The projection budget constraints need active management to prevent capacity saturation.
    \item \textbf{Physical Environment Validation:} The baseline evaluation has been restricted to simulated streams (Split-CIFAR100). Physical on-device quadruped locomotion and object recognition tests on the Unitree Go2 in real-world terrains are yet to be executed.
\end{itemize}

\section{Future Work}
\label{sec:future_work}

Building on the established baseline, the next stages of research will focus on the following milestones:

\subsection{Short-Term Implementations}
\label{subsec:short_term}
\begin{itemize}
    \item \textbf{CVAE SWR Model Training:} Finalize the CVAE generator codebase (\code{swr\_}\allowbreak\code{generator.py}) and train it using the statistics stored in the stats buffer.
    \item \textbf{Dynamic Rank Pruning:} Develop a dynamic pruning mechanism for the null-space projector to discard obsolete or redundant gradient projection directions, thereby restoring plasticity for late-stream tasks.
    \item \textbf{ImageNet-R Benchmarking:} Evaluate the complete HippoCortex pipeline on ImageNet-R to assess robustness under domain and rendition shifts.
    \item \textbf{Vision Mamba (Vim) Integration:} Integrate the bidirectional Vision Mamba (Vim) \cite{zhuVisionMambaEfficient2024} architecture as an alternative backbone options, comparing the effects of 1D causal scanning versus bidirectional 2D scanning on continual representation stability.
\end{itemize}

\subsection{Long-Term Research Directions}
\label{subsec:long_term}
\begin{itemize}
    \item \textbf{Physical Edge Deployment:} Compile the sequence backbone to TensorRT, integrate it with ROS2 and CycloneDDS on the physical Unitree Go2 quadruped, and test real-time terrain adaptation.
    \item \textbf{Federated SWR:} Explore setups where multiple quadrupeds share latent statistics during sleep to cooperatively learn across environments.
\end{itemize}

\section{Recommendations}
\label{sec:recommendations}

Based on the preliminary findings, the following recommendations are made:
\begin{itemize}
    \item \textbf{For Practitioners:} Avoid experience replay buffers of raw sensor inputs on edge robots. Focus instead on low-dimensional latent-space replay to minimize storage and memory write cycles.
    \item \textbf{For Researchers:} Prioritize dynamic basis consolidation. Null-space updates lead to rapid capacity saturation.
\end{itemize}

\section{Concluding Remarks}
\label{sec:concluding_remarks}

This research contributes to the development of sustainable, lifelong learning for mobile robots. By combining the linear-time sequence modeling of selective State Space Models with a decoupled sleep-wake generative replay paradigm, HippoCortex addresses both the stability-plasticity dilemma and the computational limits of edge deployment. The successful reproduction of the Mamba-CL baseline validates the core gradient projection mechanism and establishes a concrete performance foundation. The subsequent phases of this work will complete the implementation of HippoCortex, connecting theoretical biological memory models with physical robotic edge intelligence.